\documentclass[11pt,a4paper]{article}
\usepackage[margin=2.4cm]{geometry}
\usepackage[T1]{fontenc}
\IfFileExists{lmodern.sty}{\usepackage{lmodern}}{}
\usepackage{booktabs}
\usepackage{tabularx}
\usepackage{array}
\usepackage{enumitem}
\usepackage{fancyhdr}
\usepackage[hidelinks]{hyperref}
\usepackage{xurl}
\setlist{itemsep=2pt,topsep=4pt}
\setlength{\parskip}{4pt}
\newcolumntype{Y}{>{\raggedright\arraybackslash}X}

\pagestyle{fancy}
\fancyhf{}
\fancyhead[L]{\small Project Proposal}
\fancyhead[R]{\small TrustBazaar}
\fancyfoot[C]{\small\thepage}
\renewcommand{\headrulewidth}{0.4pt}

\begin{document}

%% ------------------------------------------------------------------ title page
\begin{titlepage}
\centering
{\large Daffodil International University\\ Department of Software Engineering\par}
\vspace{3.2cm}
{\Large Project Proposal\par}
\vspace{0.6cm}
{\LARGE\bfseries TrustBazaar: A Verified Peer-to-Peer Auction\\[4pt]
and Marketplace Platform for Used Goods\par}
\vspace{2.8cm}
{\large Submitted by\par}
\vspace{0.6cm}
\begin{tabular}{ll}
S.\ M.\ Nihal Ahmed & ID: 232-35-002\\
Afrim Hossen Khan & ID: 232-35-076\\
Sabikun Nahar Sinthia & ID: 232-35-475\\
\end{tabular}
\vspace{1.4cm}

{BSc in Software Engineering\\ Daffodil International University (DIU)\par}
\vspace{0.8cm}
{Submitted for: Capstone Project Approval\par}
\vspace{0.4cm}
{\small Supplementary document: Market Analysis of Bikroy.com and Its Alternatives\par}
\vfill
{\large October 2026\par}
\end{titlepage}

\tableofcontents
\newpage

%% ================================================================== 1 INTRODUCTION
\section{Introduction}

\subsection{Background Overview}

Second-hand goods change hands constantly in Bangladesh, but almost entirely through unmoderated
Facebook groups and general classifieds boards such as Bikroy. None of these venues solve the
problem that actually makes resale between strangers hard: nobody vouches for either side, nobody
holds the money until the item is confirmed delivered, and if a seller wants to run a genuine
auction rather than a fixed price, there is no mechanism to stop last-second sniping, fake bidding,
or a buyer who simply vanishes after ``winning.'' Delivery is negotiated ad hoc over chat, with no
fare transparency and no fallback if the courier or the buyer disappears.

The market involved is large. Bikroy, the largest classifieds board, has operated since 2012 and
carries about 297{,}000 live ads, most of them used phones, electronics and household goods,
with about 400{,}000 monthly buyers~\cite[\S2.3--2.4]{supp}. Yet its own 2014 audit found only 66\%
of sampled ads genuine, and no later audit has been published~\cite[\S2.1]{supp}.

\subsubsection*{Why Existing Platforms Don't Solve This}

It is worth being specific about why the obvious alternatives fall short, since the case for
building something new rests on it.

Facebook Marketplace and the buy-and-sell groups around it offer no identity verification beyond a
social login and no payment handling at all~\cite[\S5.2]{supp}. They also have no real auction
mechanic: ``highest offer in the comments'' has no enforced end time and leaves no record of who
actually committed to buy, so a seller who tries to auction an item has nothing to hold the winner
to.

Bikroy and similar classifieds boards are listing directories, not transaction platforms. They
connect a buyer and a seller and then step out of the way, so whatever happens with payment and
delivery happens entirely off-platform, with no escrow and no dispute mechanism. Bikroy's own terms
make this explicit: it states that it is not a party to transactions between users and caps its
liability at BDT~6{,}700~\cite[\S2.1]{supp}.

Daraz and other e-commerce marketplaces do offer payment and delivery
infrastructure~\cite[\S5.3]{supp}, but they are built for businesses selling new inventory at scale,
not for one individual selling one used item to another. There is no auction format, and the
seller-verification bar is built around business registration rather than personal KYC.

None of these platforms combine identity verification, escrow, fair auctions, and delivery
coordination into a single flow for individual resale~\cite[\S5.5]{supp}, which is precisely the
space TrustBazaar is meant to occupy.

\subsection{Problem Statement}

A buyer and seller who have never met are being asked to trust each other with real money and an
item, with no institution standing behind either promise. Four consequences follow directly from
that, and each is a concrete reason informal channels fail in practice.

The first is \textbf{no accountability}. Anyone can create a profile and list an item, and nothing
ties a listing to a real, verifiable person, so a defrauded buyer has no recourse and a scammer can
simply reappear under a new profile. Bikroy's phone-number check, for example, is defeated by
registered SIM cards bought for Tk~350--500~\cite[\S3.2]{supp}.

The second is \textbf{no payment safety}. Money changes hands either before or after delivery, with
no neutral holding point in between, so either the buyer risks paying for an item that never arrives
or the seller risks shipping an item that is never paid for. Advance payment to fake sellers is one
of the scam patterns documented on Bikroy, and the escrow rule Bangladesh Bank introduced for
e-commerce in 2021 does not reach person-to-person deals~\cite[\S3.2, \S3.5]{supp}.

The third is \textbf{no fair price discovery}. Where sellers do attempt informal ``highest offer
wins'' auctions in comment threads, there is no enforced end time, no protection against a bid
placed in the final second, and no verifiable record of who actually won, so the seller has no way to
hold the winning bidder to the deal.

The fourth is \textbf{no safe way to trade game accounts}. Accounts for games such as PUBG Mobile,
Free Fire, Clash of Clans and Steam cannot be inspected like a physical item. They are traded through
Facebook-group ``admin deals'', in which one unverified person holds both the buyer's money and the
seller's password, and even after an honest exchange the seller can recover the account after the
sale~\cite[\S5.2]{supp}.

Delivery adds a further, quieter problem. Without a system-computed price, courier cost becomes
another point of dispute layered on top of an already low-trust transaction, and there is no accepted
authority to say whether a quoted fare is reasonable. Meeting in person to avoid that cost carries its
own risk, since cash meetings have led to robbery and counterfeit-cash fraud against both buyers and
sellers~\cite[\S3.2]{supp}.

\subsection{Objectives}

\subsubsection*{Aim}

To design and build a web-based marketplace in which identity verification, escrow-held payment, a
fair time-boxed auction mechanism, and delivery coordination work as parts of one settlement flow,
for the specific case of individual resale of used goods in Bangladesh.

\subsubsection*{Specific Objectives}

\begin{enumerate}
  \item Gate the ability to list or transact behind a KYC check, so every seller and buyer in a
  completed deal is an identity-verified individual, not an anonymous profile.
  \item Build an auction mechanism whose end time and highest bid are decided by the server rather
  than the client, so a fast connection or a manipulated timestamp cannot win an auction unfairly,
  and add an automatic short extension when a bid arrives near the scheduled close, so a genuine
  last-second sniper cannot win by timing alone.
  \item Route every payment, whether by mobile financial service or card, through a single escrow
  step that only releases funds to the seller once the buyer confirms receipt, or automatically after
  a fixed timeout if the buyer does nothing.
  \item Compute the distance between buyer and seller automatically and use it either to suggest a
  safe public meetup point or to book a courier at a fare the system, not either party, has quoted,
  removing fare negotiation as a point of friction.
  \item Keep a readable, auditable record of how each payment was split between seller, courier, and
  platform, so that a dispute can be resolved by inspecting data rather than by taking someone's word
  for it.
  \item Provide an escrow flow for game-account sales that verifies the seller's control of the
  account before listing, delivers the login details so that only the buyer can read them, and holds
  the payment through a protection period after the buyer takes the account over.
\end{enumerate}

\subsection{Scope}

\subsubsection*{What the Project Covers}

The system is a single responsive web application (no native mobile app) built around one
transaction lifecycle: a user verifies their identity, lists or finds an item, buys it outright or
wins it at auction, pays into escrow, receives it via self-pickup or courier, and confirms receipt so
funds are released. Game accounts follow the same lifecycle, with delivery replaced by a protected
credential handover (Section~\ref{sec:game}). Chat, ratings, dispute handling, and notifications exist
to support that lifecycle; they are supporting infrastructure, not independent features built for
their own sake.

\subsubsection*{What It Deliberately Excludes}

The project is not attempting to be a general e-commerce storefront for new or branded retail, and it
is not building a back-office management system. The one internal-facing role in the system, used for
KYC review and dispute adjudication, is kept deliberately narrow: giving it general create, edit and
delete access over listings, users, or orders would turn the project into exactly the kind of generic
admin panel it is trying to avoid being. TrustBazaar does not integrate with game publishers and cannot
prevent a publisher from banning a traded account; this is disclosed to buyers before payment. Real
financial transactions and production courier bookings are also out of scope for the academic
deliverable; sandbox and test-mode integrations stand in for them, and that substitution is disclosed
rather than hidden behind a convincing-looking demo.

\subsection{Stakeholders}

The two primary stakeholders are the platform's users. \textbf{Buyers} are identity-verified
individuals purchasing used goods or game accounts, either at a fixed price or through the auction
mechanism. \textbf{Sellers} are identity-verified individuals listing used goods for fixed-price sale
or auction, who receive their payout once escrow releases.

The \textbf{Platform Operator} is the narrow internal-facing role responsible only for KYC document
review and dispute adjudication; the Scope section explains why this role is kept deliberately
limited.

Two groups of external providers connect to the system. The \textbf{payment providers} are bKash,
Nagad and Stripe (card, test mode), whose sandbox APIs the escrow settlement flow integrates with. The
\textbf{courier providers} fulfil courier-based deliveries through a courier sandbox API where one is
available, or through the documented mock courier service otherwise.

Finally, the \textbf{development team} (S.\ M.\ Nihal Ahmed, Afrim Hossen Khan and Sabikun Nahar
Sinthia) is responsible for design, implementation and testing, and the \textbf{academic supervisor
and evaluators} of the Department of Software Engineering, Daffodil International University, review
and assess the capstone deliverable against this proposal and the accompanying SRS.

\subsection{Proposed Solution}

TrustBazaar is a proposal to close the gap described above for one well-defined slice of the market:
individual-to-individual resale of used goods. Rather than trying to be a general marketplace, the
project deliberately narrows itself to four problems that are genuinely hard to get right (identity
trust, fair auctions, protected payment, and delivery coordination) and treats them as one connected
system instead of four separate features bolted together after the fact. Each of these answers a
specific loophole identified in the accompanying market analysis~\cite[\S3.6, \S4]{supp}.

A project like this earns its place as a capstone specifically because its interesting parts cannot be
faked with a simple database and a form. First, concurrency has to actually be handled, not assumed
away: two people bidding on the same auction in the same second, or two buyers trying to buy the last
unit of a fixed-price listing at the same moment, are ordinary situations the system has to get right
every single time, not edge cases to patch later once someone complains.

Second, multiple external payment and delivery providers have to agree with each other. A bKash
payment, a Nagad payment, and a card payment all have to collapse into the same internal notion of
``this order is paid,'' despite arriving through three different APIs with three different confirmation
mechanisms, and a courier's fare quote has to reconcile with what the platform actually charged the
buyer.

Third, trust decisions have consequences that compound. A KYC approval, a bid, and a payment release
are not independent actions: each one changes what the system is allowed to let happen next, which is a
materially different design problem from a typical CRUD workflow where most actions are reversible and
independent of each other.

By the end of the project, a demonstrable web application should exist in which a verified user can
list an item, another verified user can either buy it or win it in a fair, server-timed auction, pay
for it safely, receive it through a system-arranged method, and trust that the money only moved when it
was supposed to. The accompanying repository, System Design Document, and test suite should make the
reasoning behind each of those guarantees inspectable, not just the working demo itself, so that the
project stands up to the question ``why should anyone trust this,'' rather than only to the question
``does it work in the demo.''

\subsection{Game-Account Escrow}
\label{sec:game}

Game accounts are sold through the same escrow as physical goods, but because a password cannot be
inspected the way a phone can, the handover needs its own sequence of safeguards. Each step is designed
to close one of the ways in which today's Facebook-group admin deals fail.

The process begins before anything is sold. When a seller lists a game account, the system issues a
one-time code that the seller places in a visible field of the game profile, such as the bio or in-game
name, and then uploads a screenshot for the operator to check. This proves that the seller actually
controls the account. Once the listing is approved, it is locked to that game's player ID, so the same
account cannot be listed or sold twice.

When a buyer purchases the account, the payment goes into escrow first, and only then does the seller
hand over the login details. The handover is sealed: the buyer's browser generates an encryption key,
the seller's browser encrypts the login details with it, and the server stores only the encrypted
result. Neither the platform nor an operator can read the password, which removes the possibility that
the middleman, the role an admin plays today, misuses it.

The buyer then has 24 hours to take the account over completely by changing the password, the recovery
e-mail and phone number, removing any linked social logins, and turning on two-step verification. If
the login details do not work, the buyer reports it within that window and the order goes to a refund
review. Once the buyer has confirmed changing the recovery e-mail, however, a later claim that the
password did not work is no longer accepted, which prevents a dishonest buyer from keeping both the
account and a refund.

Finally, the payment stays in escrow for a seven-day protection period after the takeover. If the
seller uses the old e-mail or phone to recover the account during that time, the buyer raises a dispute
while the money is still held, so a refund is actually possible. After the money is released, a seller
who is found to have recovered an account is suspended and their NID is blacklisted from the platform.

Most game publishers prohibit account transfers and may ban an account that has changed
hands~\cite[\S5.2]{supp}. For that reason, every game-account listing carries a mandatory disclosure
that the buyer must acknowledge before paying, and the platform's protection is described as covering
the handover, the takeover and recovery by the seller during the protection period, not later action by
the publisher.

\subsection{AI Features}

The following AI-assisted capabilities are proposed as enhancements layered on top of the core trust
system described above, not as a replacement for any of it: the escrow, KYC, and auction mechanisms
must work correctly on their own, and AI features make specific parts of that system smarter or faster
to use. Given the project's three-month timeline, each is scoped to use existing pretrained models or
managed APIs rather than training a model from scratch, and each is treated as Should-Have or Could-Have
priority (see the Timeline and Work Plan), so the Must-Have trust guarantees are never traded off to make
room for them.

\subsubsection*{Fraud and Anomaly Detection on Bidding Patterns}

The auction engine (FR-AUC-01--FR-AUC-18) logs every bid with its bidder, amount, and timestamp, data that
already exists for the concurrency-correctness reasons described in System Requirements. This feature
adds a lightweight anomaly-scoring pass over that log to flag suspicious patterns, such as a cluster of
small proxy bids from newly created accounts driving up a price, or a bidder who consistently bids only
in the final seconds across many auctions. Flagged activity surfaces to the Platform Operator's
dispute and review queue rather than auto-blocking a bid, since a false positive here should never
silently cost a legitimate bidder an auction they fairly won.

\subsubsection*{AI-Assisted KYC Document Verification}

The KYC module (FR-KYC-01--FR-KYC-12) currently routes every submitted document to manual operator review. This
feature adds an automated pre-check step before that review, verifying that an uploaded ID photo is
legible and correctly cropped and that the selfie plausibly matches the ID photo, using an existing
document-verification or face-match API rather than a custom-trained model. This does not replace the
operator's approval; it only filters out unusable submissions (blurry photos, mismatched selfies) before
they reach the review queue, reducing manual review load without changing who has final approval
authority.

\subsubsection*{Listing Recommendation Engine}

Once enough listings and purchase history exist, a recommendation engine can surface listings a buyer is
more likely to want on the browse and search page (FR-SRC-01--FR-SRC-06), using a similarity-based approach
(category, price range, and past interaction history) rather than a deep learning model that would need
far more data than a three-month, single-market pilot can realistically generate. This is explicitly a
Could-Have: it improves discovery, but the platform is fully usable through search and filtering alone
without it.

\subsubsection*{Auto-Generated Listing Descriptions from Photos}

To lower the friction of creating a listing (FR-LST-01--FR-LST-13), a seller could optionally have a draft title
and description generated from their uploaded photos via an image-captioning or vision-language API,
which they review and edit before publishing. It is never auto-published without the seller's
confirmation, since an inaccurate AI-generated description of a used item's condition would directly
undermine the trust the rest of the platform is built around.

\subsubsection*{Scope Note}

All four features depend on the underlying data (bids, KYC documents, listings, photos) that the core
system already produces, so they are additive integrations rather than a parallel system, and each can be
descoped independently if time runs short without affecting the Must-Have deliverable described in
Section~1.6.

%% ================================================================== 2 REQUIREMENTS
\section{System Requirements}

The full requirements are specified in the accompanying Software Requirements Specification
(\texttt{docs/srs.tex}, version 2.0, IEEE 830-1998 format). It contains 191 functional requirements in
17 areas, 86 security requirements and 41 other non-functional requirements, each with a MoSCoW priority
and a verification method, together with the configurable parameters (time limits, thresholds and fees)
that the requirements refer to. Each functional area has its own identifier prefix, such as FR-AUC for the
auction engine, and maps to one backend module. This section summarizes that specification at the level
relevant to project approval, rather than reproducing it.

\subsection{Functional Requirement Areas}

The requirements follow the transaction lifecycle. A user first registers and signs in (FR-ACC) and
passes KYC verification (FR-KYC), after which they can create fixed-price or auction listings (FR-LST)
and find items through search and discovery (FR-SRC). Purchases then take one of two paths: the auction
engine, with server-authoritative bidding, proxy bids, a reserve price and anti-sniping (FR-AUC), or
fixed-price purchase with a concurrency-safe stock reservation (FR-FIX).

Every purchase is settled through payment and escrow (FR-PAY), including seller payouts and refunds,
and fulfilled either through the delivery area (FR-DLV), which chooses between a safe meetup with a
handover code and a courier, or, for game accounts, through the game-account escrow (FR-GAM), which adds
ownership proof, sealed credential handover, guided takeover and the protection hold described in
Section~\ref{sec:game}. Around this core sit the supporting areas: real-time chat (FR-CHT); ratings,
reviews and trust signals (FR-REV); dispute resolution (FR-DSP); and notifications (FR-NTF).
Monetization, covering the platform commission and delivery margin, is specified in FR-MON;
trust and safety controls in FR-TNS; the deliberately narrow operator and administrator functions in
FR-OPS; and privacy and data-subject rights in FR-PRV. Security requirements (SEC) are collected in a
separate chapter of the SRS, based on OWASP ASVS 5.0 Level~2.

\subsection{Key Non-Functional Requirements}

The requirements that most directly shape the architecture are the ones concerning correctness under
concurrency and real-time state propagation. The auction engine's end time and highest bid must be
authoritative on the server regardless of client-side timestamps or connection speed; live bid and chat
updates must propagate within approximately one second; and the system must run as separate,
independently restartable containers (the backend API with its real-time gateway, the web client, the
database and the cache), so that no single process holding in-memory state can silently corrupt an
in-progress auction.

\subsection{Acceptance Criteria}

Rather than judging the project by feature count, the following properties are treated as the actual
deliverable, since the rest of the system exists to give them a context to operate in:

\begin{enumerate}
  \item Two simultaneous bids on the same auction never both register as the winning bid.
  \item Two simultaneous buyers can never both purchase the last unit of a fixed-price listing.
  \item An escrow payment is never released to a seller without either explicit buyer confirmation or a
  genuine timeout.
  \item A payment webhook delivered twice by a provider never results in the order being charged or
  settled twice.
  \item Game-account login details are never stored on the server in a form the server or an operator
  can read.
  \item Payment for a game-account order is never released before the protection period ends, unless the
  buyer releases it early or an operator decides a dispute.
\end{enumerate}

These are verified by an automated test suite that specifically targets simultaneous bidding,
simultaneous stock decrement, repeated or duplicate payment-webhook delivery, the correctness of the
fee-split arithmetic, and the game-account release rules. The suite runs automatically on every push and
pull request to the main branch, so a regression is caught before it ever reaches a demonstration.

%% ================================================================== 3 TOOLS
\section{Tools and Technologies}

The technology stack below is the one set up in the project repository~\cite{repo}, which already
contains the module structure, data model, container setup and continuous-integration pipeline that
the implementation will build on. Each choice is made for a reason specific to this project rather than
by default.

\subsection{Backend}

The backend is a \textbf{NestJS 10} application written in \textbf{TypeScript 5} and running on
\textbf{Node.js 20}. NestJS was chosen because its module system maps directly onto the SRS: the
repository has one module per feature area (authentication, users, KYC, listings, auctions, orders,
payments, delivery, chat, reviews, disputes and notifications), so each requirement range has an obvious
home and the three team members can work on separate modules without stepping on each other. Every
incoming request passes through a global validation layer built on \textbf{class-validator} and
\textbf{class-transformer}, which checks the request against its expected shape and rejects any fields the
endpoint does not expect, so malformed or over-posted data never reaches the business logic.

Users will authenticate with \textbf{JSON Web Tokens}, issued through \textbf{@nestjs/jwt} and checked
by \textbf{Passport} with the \textbf{passport-jwt} strategy, and passwords will be stored only as
\textbf{Argon2id} hashes, which replace the \textbf{bcrypt} library currently in the project because bcrypt
ignores input beyond 72 bytes and would truncate long Bangla passphrases. Session tokens are carried in
\texttt{HttpOnly} cookies rather than in browser storage, so injected scripts cannot read them. Configuration such as secrets, database addresses and provider keys is loaded from
environment variables through \textbf{@nestjs/config}, so no credential is ever written into the code.

\subsection{Data Storage}

\textbf{PostgreSQL 15} is chosen over a document store because settlement and bidding correctness
genuinely depend on transactional (ACID) guarantees; this is a case where that property is
load-bearing, not a formality that could be swapped out later. The application reaches it through the
\textbf{TypeORM 0.3} object-relational mapper and the \textbf{pg} driver. The data model is already
defined as fourteen entities covering users, KYC documents, listings, bids, orders, escrow transactions,
the settlement ledger, delivery requests, chat threads and messages, reviews, disputes, featured listings
and subscriptions. The escrow transaction records its state (held, released, refunded, or frozen under
dispute) together with an append-only history of every change, which is what allows a dispute to be
decided by inspecting the record rather than by taking either party's word. Automatic schema
synchronisation is used only during early development and will be replaced by versioned migrations before
any real data exists.

\textbf{Redis 7}, accessed through the \textbf{ioredis} client, holds auction countdown state and
coordinates bid locks outside any single server process's memory, so restarting or scaling the backend
cannot accidentally reset or extend a live auction's true end time. The same Redis instance provides
publish/subscribe messaging between server instances and the counters used for rate limiting.

\subsection{Real-Time Communication}

Live bid updates, countdown synchronisation, chat delivery and notifications are carried over
\textbf{Socket.IO 4} through a NestJS WebSocket gateway, since a marketplace where the displayed highest bid
is stale by several seconds undermines the fairness the auction engine is meant to guarantee in the first
place. Each auction and each chat thread will be its own room, so an accepted bid is broadcast only to the
people watching that listing; the gateway and its bid-broadcast method are already in place. When more than one backend instance is running, the gateway relies on Redis
publish/subscribe rather than keeping connection state in a single process's memory.

\subsection{Frontend}

The web client is built with \textbf{React 18} and \textbf{TypeScript}, using \textbf{React Router 6} for
navigation and \textbf{Vite 7} as the development server and build tool. The client talks to the backend
through a small API layer that sends each request with the session cookie set by the backend. The first screens set up in the
repository are the home page, login, the KYC submission wizard and the listing detail page, which together
cover the entry point of the transaction lifecycle. For the game-account escrow, the browser's built-in
\textbf{Web Crypto API} will perform the key generation and encryption of login details on the client, so
the server never handles a readable password and no custom cryptography is written.

\subsection{External Integrations}

Payments use each provider's sandbox or test mode: the \textbf{bKash} tokenized-checkout sandbox, the
\textbf{Nagad} merchant sandbox, and \textbf{Stripe} in test mode for cards, with Stripe's webhook signing
secret used to verify that payment confirmations genuinely come from the provider. For delivery, a single
configuration setting selects the courier provider: a documented \textbf{mock courier} by default, or the
\textbf{Pathao} or \textbf{Steadfast} courier APIs where sandbox access is available. Where no public courier
sandbox exists, the mock service is used and clearly labelled as such rather than presented as a live
integration.

\subsection{Infrastructure, Testing and Quality}

\textbf{Docker Compose} runs the system as four separate containers (PostgreSQL, Redis, the backend and the
frontend) built on lightweight Alpine Linux images. Health checks make the backend wait until the database
and cache are actually ready before it starts, and the arrangement mirrors how the pieces would be deployed,
keeping local development close to that reality rather than diverging from it.

Tests are written with \textbf{Jest 29} and \textbf{ts-jest}. Alongside ordinary unit tests, the repository
already defines a dedicated concurrency test suite whose cases correspond to the acceptance criteria:
simultaneous bidding, the anti-sniping extension, proxy bidding, the last-unit purchase race, duplicate
payment webhooks and fee-split reconciliation. Each case is currently marked as pending and is filled in as
the feature behind it is built, and a feature is not considered finished until its case passes. Code style and common errors are checked by \textbf{ESLint} with the
TypeScript rules.

Every push and pull request to the main branch triggers a \textbf{GitHub Actions} pipeline that lints the
backend, runs the unit tests, and runs the concurrency suite against real PostgreSQL and Redis service
containers, then builds the frontend, so a regression is caught before it is merged. \textbf{GitHub CodeQL}
scans the TypeScript code for security weaknesses on every change and again weekly, and
\textbf{Dependabot} proposes weekly minor and patch updates to dependencies, while major-version upgrades are
held back for deliberate, manual handling.

%% ================================================================== 4 PLAN
\section{Project Timeline and Work Plan}

\subsection{Development Approach}

The build proceeds as a sequence of working, demonstrable increments rather than a single delivery at
the end. Each increment adds one complete slice of the transaction lifecycle (data model, backend logic,
and the corresponding UI together), so that at every stage there is something that can actually be
clicked through, not just a partially wired backend sitting behind a placeholder screen.

\subsubsection*{Planned Increments}

The first increment establishes identity and accounts: registration, login, and the KYC
submission-and-review path that gates everything downstream. The second adds listings, covering the
creation of fixed-price and auction listings and the browsing, searching and filtering of the catalogue
by category, price, condition and distance.

The third and fourth increments build the two purchase paths. The auction path brings in
server-authoritative bidding, automatic proxy bids, the reserve price and the anti-sniping extension,
tested specifically under simultaneous bidders rather than only one bidder at a time. The fixed-price path
adds purchase with a stock decrement that stays correct when several buyers try to claim the last unit at
once.

The fifth increment connects payment and escrow, integrating the mobile financial services and card
processor sandboxes with the hold-then-release settlement logic behind both. The sixth adds delivery: the
distance computation and the self-pickup or courier decision, using a courier sandbox or its documented mock
substitute. The seventh builds the trust and communication layer of chat, ratings, disputes and
notifications, wired to the specific transactions they support rather than built as generic messaging.

The eighth increment adds the game-account escrow (ownership proof, sealed credentials, guided takeover and
the protection hold), built on the escrow and dispute logic already in place. The ninth and last covers
monetization and hardening: the commission and delivery-margin calculation on top of the existing ledger,
followed by a pass on rate limiting, audit logging, and behaviour under concurrent load.

\subsection{Team and Roles}

The project is undertaken by a three-member team: S.\ M.\ Nihal Ahmed (ID: 232-35-002), Afrim Hossen Khan
(ID: 232-35-076) and Sabikun Nahar Sinthia (ID: 232-35-475).

Having three people changes the methodology in one important way: instead of the strictly sequential,
one-slice-at-a-time build described for a solo effort, the workstreams below are split across team members
and run partly in parallel once the shared foundation (data model and authentication) is in place, which is
what makes a three-month timeline realistic rather than optimistic. A shared Git repository with feature
branches and a lightweight code-review step before merging to main keeps the three workstreams from silently
drifting apart. Development, sandbox accounts, and hosting are expected to run on free or student-tier
offerings (payment provider sandboxes, a free-tier cloud instance or local Docker Compose setup for
demonstration), keeping the resource footprint of the build close to zero cost.

\subsection{Timeline (12 Weeks / 3 Months)}

The three core workstreams (the auction and listings path, the payment and escrow path, and the delivery and
trust-layer path) are assigned to the three team members and developed in parallel from Week~3 onward,
against the shared data model and authentication built together in Weeks~1--2.

\begin{center}
\small
\begin{tabularx}{\textwidth}{@{}l Y l@{}}
  \toprule
  \textbf{Weeks} & \textbf{Focus} & \textbf{Owner(s)} \\
  \midrule
  1--2 & Shared setup: repo, Docker Compose, data model, auth, KYC flow & All three, together \\
  3--5 & Auction engine, fixed-price purchase, concurrency tests & Member A \\
  3--5 & Payment/escrow integration, settlement ledger & Member B \\
  3--5 & Listings, search/discovery, delivery decision engine & Member C \\
  6--7 & Integrate the three workstreams into one order lifecycle & All three, together \\
  8 & Chat and notifications & Member A \\
  8 & Game-account escrow: sealed credentials, ownership proof & Member B \\
  8 & Ratings and disputes & Member C \\
  9 & Monetization features (commission, delivery margin) & Member A \\
  9 & Game-account escrow: guided takeover, protection hold & Member B \\
  9 & Security hardening: rate-limiting, audit logging, IDOR checks & Member C, reviewed by all \\
  10 & End-to-end and load testing of concurrency-sensitive paths & All three, together \\
  11 & Bug fixing, System Design Document, README & All three, together \\
  12 & Final polish, demonstration rehearsal, submission & All three, together \\
  \bottomrule
\end{tabularx}
\end{center}

Member assignment (A/B/C) is deliberately left as a role rather than pinned to a specific name here, so the
team can assign it based on who ends up strongest in each area during Weeks~1--2. The auction engine and the
payment and escrow flow remain the critical path, since chat, ratings, disputes, monetization and the
game-account escrow all depend on an order existing and having been paid for correctly, so those two
workstreams get first claim on time if Weeks~3--5 run long. If the schedule slips, the features scoped down
first are the Should-Have and Could-Have ones (delivery automation, notifications, featured listings, the Pro
Seller tier), since the Must-Have core (KYC, the auction engine, escrow settlement, chat, disputes and the
game-account escrow) is what the project's central claim depends on.

\subsection{Risk Management}

The project carries five main risks, each with a planned response.

\textbf{Sandbox access delays.} Developer sandbox approval for bKash and Nagad can be slow to obtain. If it
is not available in time, Stripe test mode alone will be used, and the limitation will be stated openly in
the final documentation rather than worked around silently.

\textbf{No public courier sandbox.} If neither Pathao nor Steadfast sandbox access is available, the mock
courier service, with realistic fare and status behaviour, stands in for them and is documented explicitly
as a substitution rather than presented as if it were live.

\textbf{Concurrency defects.} This is the single biggest technical risk. It is mitigated by row-level
locking in PostgreSQL and by writing the concurrency tests before the corresponding feature is considered
finished, not after a bug is reported; the cases of the concurrency suite are already defined in the
repository and run in continuous integration on every change.

\textbf{Publisher terms on game accounts.} Publishers may ban accounts that have been
traded~\cite[\S5.2]{supp}. The response is threefold: buyers must acknowledge a mandatory disclosure before
paying, the platform's protection is limited to the handover and the protection period, and the
game-account listing type sits behind a configuration switch, so it can be disabled without affecting the
rest of the platform if that becomes necessary.

\textbf{Account recovery after release.} An original owner may still recover a sold account through the
publisher's own support process, which no marketplace can fully prevent. The seven-day protection hold
reduces this risk, and identity-based accountability for the seller deals with it if it happens after the
money has been released.

%% ================================================================== REFERENCES
\begin{thebibliography}{9}
\small
\bibitem{supp} S.~M.~N.~Ahmed, A.~H.~Khan and S.~N.~Sinthia, ``Trust in Bangladesh's Second-Hand
Online Market: A Market Analysis of Bikroy.com and Its Alternatives,'' supplementary report to the
TrustBazaar project proposal, Department of Software Engineering, Daffodil International University,
October 2026.
\bibitem{repo} TrustBazaar project repository. \url{https://github.com/nihal4/trustbazaar}
(accessed 7 Oct 2026).
\end{thebibliography}

\end{document}
